\section{Exercise 1}

Trước tiên, chương trình khai báo các chuỗi thông báo và các hằng số thực cần thiết trong vùng nhớ \texttt{.data}:

\begin{codeblock}[language=MIPS]
.data
  prompt_a: .asciiz "Please insert a: "
  prompt_b: .asciiz "Please insert b: "
  prompt_c: .asciiz "Please insert c: "

  msg_x1:   .asciiz "x1 = "
  msg_x2:   .asciiz " and x2 = "
  msg_one:  .asciiz "There is one solution, x = "
  msg_none: .asciiz "There is no real solution"
  msg_inf:  .asciiz "There are infinite real solutions"

  const0f:  .float 0.0
  const2f:  .float 2.0
  const4f:  .float 4.0
\end{codeblock}

Tiếp theo, chương trình lần lượt đọc ba hệ số $a, b, c$ từ bàn phím và lưu vào các thanh ghi \texttt{\$f20}, \texttt{\$f21}, \texttt{\$f22}. Khối lệnh nhập $b$ và $c$ có cấu trúc tương tự như khi nhập $a$:

\begin{codeblock}[language=MIPS]
  li $v0, 4        # syscall 4: in chuỗi
  la $a0, prompt_a # in "Please insert a: "
  syscall

  li $v0, 6        # syscall 6: đọc số thực
  syscall
  mov.s $f20, $f0  # $f20 = a (lưu lại trước khi $f0 bị ghi đè)
\end{codeblock}

Trước khi tính $\Delta$, chương trình kiểm tra điều kiện $a = 0$. Nếu $a = 0$, phương trình suy biến thành dạng bậc nhất $bx + c = 0$ và được rẽ nhánh để xử lý riêng:

\begin{codeblock}[language=MIPS]
  l.s $f0, const0f # $f0 = 0.0 làm mốc so sánh

  c.eq.s $f20, $f0 # so sánh a với 0.0
  bc1t linear_case # a == 0 thì rẽ sang nhánh bậc nhất
\end{codeblock}

Với $a \neq 0$, chương trình tính $\Delta = b^2 - 4ac$ và đối chiếu với $0$ để phân loại số nghiệm ($\Delta < 0$: vô nghiệm; $\Delta = 0$: nghiệm kép; $\Delta > 0$: hai nghiệm phân biệt):

\begin{codeblock}[language=MIPS]
  mul.s $f1, $f21, $f21 # $f1 = b^2
  l.s $f2, const4f      # $f2 = 4.0
  mul.s $f3, $f20, $f22 # $f3 = a * c
  mul.s $f3, $f3, $f2   # $f3 = 4ac
  sub.s $f3, $f1, $f3   # $f3 = delta

  c.lt.s $f3, $f0       # delta < 0
  bc1t no_solution      # rẽ nhánh vô nghiệm

  c.eq.s $f3, $f0       # delta == 0
  bc1t one_solution     # rẽ nhánh nghiệm kép
\end{codeblock}

Khi $\Delta > 0$, hai nghiệm phân biệt được tính theo công thức $x_{1,2} = \dfrac{-b \pm \sqrt{\Delta}}{2a}$:

\begin{codeblock}[language=MIPS]
  sqrt.s $f4, $f3      # $f4 = sqrt(delta)
  neg.s $f5, $f21      # $f5 = -b
  l.s $f6, const2f     # $f6 = 2.0
  mul.s $f6, $f6, $f20 # $f6 = 2a

  sub.s $f7, $f5, $f4  # $f7 = -b - sqrt(delta)
  div.s $f8, $f7, $f6  # $f8 = x1

  add.s $f9, $f5, $f4  # $f9 = -b + sqrt(delta)
  div.s $f10, $f9, $f6 # $f10 = x2
\end{codeblock}

Hai nghiệm sau đó được in ra màn hình thông qua \texttt{syscall 2} (giá trị cần in được nạp vào \texttt{\$f12}), rồi kết thúc chương trình:

\begin{codeblock}[language=MIPS]
  li $v0, 4        # syscall 4: in chuỗi
  la $a0, msg_x1   # in "x1 = "
  syscall

  li $v0, 2        # syscall 2: in số thực
  mov.s $f12, $f8  # $f12 = x1
  syscall

  li $v0, 4        # syscall 4: in chuỗi
  la $a0, msg_x2   # in " and x2 = "
  syscall

  li $v0, 2        # syscall 2: in số thực
  mov.s $f12, $f10 # $f12 = x2
  syscall

  j exit           # kết thúc chương trình sau khi in 2 nghiệm
\end{codeblock}

Ở nhánh $a = 0$, nếu $b \neq 0$ thì phương trình bậc nhất có nghiệm duy nhất $x = -\dfrac{c}{b}$:

\begin{codeblock}[language=MIPS]
linear_case:           # nhánh a == 0: bx + c = 0
  c.eq.s $f21, $f0     # so sánh b với 0.0
  bc1t check_c         # b == 0 thì xét tiếp c

  neg.s $f5, $f22      # $f5 = -c
  div.s $f6, $f5, $f21 # $f6 = -c / b

  li $v0, 4            # syscall 4: in chuỗi
  la $a0, msg_one      # in "There is one solution, x = "
  syscall

  li $v0, 2            # syscall 2: in số thực
  mov.s $f12, $f6      # $f12 = -c / b
  syscall

  j exit
\end{codeblock}

Khi $\Delta = 0$, phương trình bậc hai có nghiệm kép $x = -\dfrac{b}{2a}$:

\begin{codeblock}[language=MIPS]
one_solution:          # delta == 0: nghiệm kép -b / 2a
  neg.s $f5, $f21      # $f5 = -b
  l.s $f6, const2f     # $f6 = 2.0
  mul.s $f6, $f6, $f20 # $f6 = 2a
  div.s $f7, $f5, $f6  # $f7 = -b / 2a

  li $v0, 4            # syscall 4: in chuỗi
  la $a0, msg_one      # in "There is one solution, x = "
  syscall

  li $v0, 2            # syscall 2: in số thực
  mov.s $f12, $f7      # $f12 = -b / 2a
  syscall

  j exit
\end{codeblock}

Khi $a = b = 0$, chương trình kiểm tra tiếp hệ số $c$: nếu $c = 0$ thì phương trình nghiệm đúng với mọi $x \in \mathbb{R}$, ngược lại ($c \neq 0$) thì vô nghiệm:

\begin{codeblock}[language=MIPS]
check_c:           # nhánh a == 0 và b == 0: xét phương trình c = 0
  c.eq.s $f22, $f0 # so sánh c với 0.0
  bc1f no_solution # c != 0 thì vô nghiệm

  li $v0, 4        # syscall 4: in chuỗi
  la $a0, msg_inf  # in "There are infinite real solutions"
  syscall

  j exit
\end{codeblock}

Label \texttt{no\_solution} xử lý chung cho cả hai trường hợp vô nghiệm ($\Delta < 0$ hoặc $a = b = 0, c \neq 0$) trước khi thoát chương trình:

\begin{codeblock}[language=MIPS]
no_solution:       # delta < 0 hoặc (a == 0, b == 0, c != 0)
  li $v0, 4        # syscall 4: in chuỗi
  la $a0, msg_none # in "There is no real solution"
  syscall

exit:              # điểm kết thúc chung
  li $v0, 10       # syscall 10: kết thúc chương trình
  syscall
\end{codeblock}

Kết quả chạy thử nghiệm với một số testcase:

\begin{codeblock}[language={}]
Please insert a: 1.2
Please insert b: 9.6
Please insert c: 1
x1 = -7.89444 and x2 = -0.10555982
\end{codeblock}

\begin{codeblock}[language={}]
Please insert a: 2.4
Please insert b: -3.2
Please insert c: 3.1
There is no real solution
\end{codeblock}

\begin{codeblock}[language={}]
Please insert a: 1
Please insert b: 2
Please insert c: 1
There is one solution, x = -1.0
\end{codeblock}

\begin{codeblock}[language={}]
Please insert a: 0
Please insert b: 0
Please insert c: 0
There are infinite real solutions
\end{codeblock}
